\section*{Capítulo I, Problema VIII}

\textbf{Problema:}

8. Sea $A$ un vector perpendicular a todo vector $X$. Demuestra que $A = 0$.

\textbf{Solución:}

Si $A$ es perpendicular a todo vector $X$, entonces el producto escalar $A \cdot X = 0$ para cualquier vector $X$. Supongamos que $A \neq 0$. Esto significa que $A$ tiene al menos un componente no nulo. Sin pérdida de generalidad, supongamos que el primer componente de $A$ es distinto de cero, es decir, $A = (a_1, a_2, \dots, a_n)$ con $a_1 \neq 0$.

Consideremos el vector $X = (1, 0, 0, \dots, 0)$. Entonces:
\[
    A \cdot X = a_1(1) + a_2(0) + \dots + a_n(0) = a_1.
\]

Dado que $A \cdot X = 0$ por hipótesis, se tiene que $a_1 = 0$. Esto contradice nuestra suposición de que $a_1 \neq 0$. Por lo tanto, todos los componentes de $A$ deben ser cero, lo que implica que $A = 0$.

Por lo tanto, hemos demostrado que si $A$ es perpendicular a todo vector $X$, entonces $A = 0$. 
